\documentclass[10pt,a4paper]{article}
\usepackage[margin=18mm]{geometry}
\usepackage{fontspec}
\usepackage{amsmath,amssymb}
\usepackage{booktabs}
\usepackage{tabularx}
\usepackage{enumitem}
\usepackage{xcolor}
\usepackage{hyperref}
\usepackage{listings}

\setmainfont{Noto Sans CJK KR}[AutoFakeSlant=0.2]
\setsansfont{Noto Sans CJK KR}
\setmonofont{D2Coding}
\setlength{\parindent}{0pt}
\setlength{\parskip}{5pt}
\setlist{nosep,leftmargin=6mm}
\pagestyle{plain}
\hypersetup{
  colorlinks=true,
  linkcolor=blue,
  urlcolor=blue,
  pdftitle={1번 암호분석: Caesar/Vigenère 복원과 학습 분류기},
  pdfauthor={암호분석경진대회 제출 보고서}
}
\lstset{
  basicstyle=\ttfamily\small,
  frame=single,
  columns=fullflexible,
  breaklines=true,
  keepspaces=true
}

\title{1번 암호분석\\Caesar/Vigenère 복원과 학습 분류기}
\author{암호분석경진대회 제출 보고서}
\date{2026년 7월}

\begin{document}
\maketitle

\section*{제출 답 요약}

\begin{center}
\begin{tabularx}{\textwidth}{lX}
\toprule
항목 & 결과 \\
\midrule
\texttt{ciphertexts1.txt} & Caesar, IC $0.067638\to\mathbf{0.068}$,
shift $6$ \\
\texttt{ciphertexts2.txt} & Vigenère, IC $0.043972\to\mathbf{0.044}$,
key \texttt{KLVOJ}, 줄마다 key position 초기화 \\
3-다/4-다 평문 &
\texttt{THESOVEREIGNTYOFTHEREPUBLICOFKOREA}\newline
\texttt{SHALLRESIDEINTHEPEOPLEANDALLSTATEAUTHORITY}\newline
\texttt{SHALLEMANATEFROMTHEPEOPLE} \\
분류기 holdout &
동일 corpus·동일 key 안의 paired line holdout $58/58$ \\
관측 가능한 1--4 구조 &
모두 period-1 Caesar-like, shift $6,10,10,6$ \\
출제 의도 추정 라벨 &
\textbf{Caesar, Vigenère, Vigenère, Caesar} \\
\bottomrule
\end{tabularx}
\end{center}

\section{Caesar와 Vigenère}

알파벳을 $A=0,\ldots,Z=25$로 나타내면 Caesar는 하나의 shift $k$를
모든 위치에 적용한다.
\[
 C_i=P_i+k\pmod {26},\qquad P_i=C_i-k\pmod {26}.
\]
따라서 영어의 불균등한 단문자 빈도와 IC가 보존되고, 빈도표의 문자
이름만 이동한다. key 공간도 26개다.

Vigenère는 길이 $r$인 key를 반복한다.
\[
 C_i=P_i+K_{i\bmod r}\pmod {26}.
\]
서로 다른 shift의 빈도분포가 섞여 전체 IC는 균등분포의
$1/26\simeq0.0385$ 쪽으로 낮아진다. 반대로 같은 $i\bmod r$ 위치끼리
모으면 각 열이 다시 Caesar가 된다. 길이 1 Vigenère는 Caesar와
수학적으로 완전히 같다는 포함 관계가 마지막 분류 문항에서 중요하다.

\section{IC 계산과 암호방식 특정}

문자별 횟수를 $n_i$, 전체 길이를 $N$이라 하면
\[
 \operatorname{IC}=\frac{\sum_i n_i(n_i-1)}{N(N-1)}
\]
를 계산했다. A--Z 이외 문자를 제거한 두 파일은 각각 19,897자다.

\begin{center}
\begin{tabular}{lrrl}
\toprule
파일 & IC & 요구 반올림 & 판정 \\
\midrule
\texttt{ciphertexts1.txt} & 0.067638 & 0.068 & Caesar \\
\texttt{ciphertexts2.txt} & 0.043972 & 0.044 & Vigenère \\
\bottomrule
\end{tabular}
\end{center}

영어 평문의 IC는 약 $0.065\text{--}0.067$로 알려져 있다(Friedman).
\texttt{ciphertexts1.txt}의 $0.068$은 이 값과 거의 같아 빈도 분포가
보존된 Caesar로, \texttt{ciphertexts2.txt}의 $0.044$는 균등분포값
$0.0385$ 쪽에 더 가까워 여러 shift가 섞인 Vigenère로 판정했다.

\section{Ceaser 평문 복원}

시저 암호로 판단한 \texttt{ciphertexts1.txt}를 가능한 key $k=0,\ldots,25$
전체로 복호화하고, 각 복호문에서 A--Z 출현 횟수를 세었다.

\begin{center}
\tiny
\setlength{\tabcolsep}{1.3pt}
\begin{tabular}{r*{26}{r}}
\toprule
$k$ & A & B & C & D & E & F & G & H & I & J & K & L & M & N & O & P & Q & R & S & T & U & V & W & X & Y & Z \\
\midrule
0 & 428 & 181 & 134 & 45 & 355 & 57 & 1676 & 459 & 751 & 734 & 2524 & 495 & 245 & 899 & 1484 & 37 & 36 & 1016 & 499 & 1427 & 1479 & 444 & 28 & 1198 & 1361 & 1905 \\
1 & 181 & 134 & 45 & 355 & 57 & 1676 & 459 & 751 & 734 & 2524 & 495 & 245 & 899 & 1484 & 37 & 36 & 1016 & 499 & 1427 & 1479 & 444 & 28 & 1198 & 1361 & 1905 & 428 \\
2 & 134 & 45 & 355 & 57 & 1676 & 459 & 751 & 734 & 2524 & 495 & 245 & 899 & 1484 & 37 & 36 & 1016 & 499 & 1427 & 1479 & 444 & 28 & 1198 & 1361 & 1905 & 428 & 181 \\
3 & 45 & 355 & 57 & 1676 & 459 & 751 & 734 & 2524 & 495 & 245 & 899 & 1484 & 37 & 36 & 1016 & 499 & 1427 & 1479 & 444 & 28 & 1198 & 1361 & 1905 & 428 & 181 & 134 \\
4 & 355 & 57 & 1676 & 459 & 751 & 734 & 2524 & 495 & 245 & 899 & 1484 & 37 & 36 & 1016 & 499 & 1427 & 1479 & 444 & 28 & 1198 & 1361 & 1905 & 428 & 181 & 134 & 45 \\
5 & 57 & 1676 & 459 & 751 & 734 & 2524 & 495 & 245 & 899 & 1484 & 37 & 36 & 1016 & 499 & 1427 & 1479 & 444 & 28 & 1198 & 1361 & 1905 & 428 & 181 & 134 & 45 & 355 \\
6 & 1676 & 459 & 751 &734 & 2524 & 495 & 245 & 899 & 1484 & 37 & 36 & 1016 & 499 & 1427 & 1479 & 444 & 28 & 1198 & 1361 & 1905 & 428 & 181 & 134 & 45 & 355 & 57 \\
7 & 459 & 751 & 734 & 2524 & 495 & 245 & 899 & 1484 & 37 & 36 & 1016 & 499 & 1427 & 1479 & 444 & 28 & 1198 & 1361 & 1905 & 428 & 181 & 134 & 45 & 355 & 57 & 1676 \\
8 & 751 & 734 & 2524 & 495 & 245 & 899 & 1484 & 37 & 36 & 1016 & 499 & 1427 & 1479 & 444 & 28 & 1198 & 1361 & 1905 & 428 & 181 & 134 & 45 & 355 & 57 & 1676 & 459 \\
9 & 734 & 2524 & 495 & 245 & 899 & 1484 & 37 & 36 & 1016 & 499 & 1427 & 1479 & 444 & 28 & 1198 & 1361 & 1905 & 428 & 181 & 134 & 45 & 355 & 57 & 1676 & 459 & 751 \\
10 & 2524 & 495 & 245 & 899 & 1484 & 37 & 36 & 1016 & 499 & 1427 & 1479 & 444 & 28 & 1198 & 1361 & 1905 & 428 & 181 & 134 & 45 & 355 & 57 & 1676 & 459 & 751 & 734 \\
11 & 495 & 245 & 899 & 1484 & 37 & 36 & 1016 & 499 & 1427 & 1479 & 444 & 28 & 1198 & 1361 & 1905 & 428 & 181 & 134 & 45 & 355 & 57 & 1676 & 459 & 751 & 734 & 2524 \\
12 & 245 & 899 & 1484 & 37 & 36 & 1016 & 499 & 1427 & 1479 & 444 & 28 & 1198 & 1361 & 1905 & 428 & 181 & 134 & 45 & 355 & 57 & 1676 & 459 & 751 & 734 & 2524 & 495 \\
13 & 899 & 1484 & 37 & 36 & 1016 & 499 & 1427 & 1479 & 444 & 28 & 1198 & 1361 & 1905 & 428 & 181 & 134 & 45 & 355 & 57 & 1676 & 459 & 751 & 734 & 2524 & 495 & 245 \\
14 & 1484 & 37 & 36 & 1016 & 499 & 1427 & 1479 & 444 & 28 & 1198 & 1361 & 1905 & 428 & 181 & 134 & 45 & 355 & 57 & 1676 & 459 & 751 & 734 & 2524 & 495 & 245 & 899 \\
15 & 37 & 36 & 1016 & 499 & 1427 & 1479 & 444 & 28 & 1198 & 1361 & 1905 & 428 & 181 & 134 & 45 & 355 & 57 & 1676 & 459 & 751 & 734 & 2524 & 495 & 245 & 899 & 1484 \\
16 & 36 & 1016 & 499 & 1427 & 1479 & 444 & 28 & 1198 & 1361 & 1905 & 428 & 181 & 134 & 45 & 355 & 57 & 1676 & 459 & 751 & 734 & 2524 & 495 & 245 & 899 & 1484 & 37 \\
17 & 1016 & 499 & 1427 & 1479 & 444 & 28 & 1198 & 1361 & 1905 & 428 & 181 & 134 & 45 & 355 & 57 & 1676 & 459 & 751 & 734 & 2524 & 495 & 245 & 899 & 1484 & 37 & 36 \\
18 & 499 & 1427 & 1479 & 444 & 28 & 1198 & 1361 & 1905 & 428 & 181 & 134 & 45 & 355 & 57 & 1676 & 459 & 751 & 734 & 2524 & 495 & 245 & 899 & 1484 & 37 & 36 & 1016 \\
19 & 1427 & 1479 & 444 & 28 & 1198 & 1361 & 1905 & 428 & 181 & 134 & 45 & 355 & 57 & 1676 & 459 & 751 & 734 & 2524 & 495 & 245 & 899 & 1484 & 37 & 36 & 1016 & 499 \\
20 & 1479 & 444 & 28 & 1198 & 1361 & 1905 & 428 & 181 & 134 & 45 & 355 & 57 & 1676 & 459 & 751 & 734 & 2524 & 495 & 245 & 899 & 1484 & 37 & 36 & 1016 & 499 & 1427 \\
21 & 444 & 28 & 1198 & 1361 & 1905 & 428 & 181 & 134 & 45 & 355 & 57 & 1676 & 459 & 751 & 734 & 2524 & 495 & 245 & 899 & 1484 & 37 & 36 & 1016 & 499 & 1427 & 1479 \\
22 & 28 & 1198 & 1361 & 1905 & 428 & 181 & 134 & 45 & 355 & 57 & 1676 & 459 & 751 & 734 & 2524 & 495 & 245 & 899 & 1484 & 37 & 36 & 1016 & 499 & 1427 & 1479 & 444 \\
23 & 1198 & 1361 & 1905 & 428 & 181 & 134 & 45 & 355 & 57 & 1676 & 459 & 751 & 734 & 2524 & 495 & 245 & 899 & 1484 & 37 & 36 & 1016 & 499 & 1427 & 1479 & 444 & 28 \\
24 & 1361 & 1905 & 428 & 181 & 134 & 45 & 355 & 57 & 1676 & 459 & 751 & 734 & 2524 & 495 & 245 & 899 & 1484 & 37 & 36 & 1016 & 499 & 1427 & 1479 & 444 & 28 & 1198 \\
25 & 1905 & 428 & 181 & 134 & 45 & 355 & 57 & 1676 & 459 & 751 & 734 & 2524 & 495 & 245 & 899 & 1484 & 37 & 36 & 1016 & 499 & 1427 & 1479 & 444 & 28 & 1198 & 1361 \\
\bottomrule
\end{tabular}
\end{center}

문제의 영어 기대빈도와 Pearson $\chi^2$ 값을 비교했다.
\[
 \chi^2(k)=\sum_i\frac{(O_i-E_i)^2}{E_i}.
\]

\begin{center}
\small
\begin{tabular}{rr@{\hspace{9mm}}rr}
\toprule
$k$ & $\chi^2$ & $k$ & $\chi^2$ \\
\midrule
0 & 370761.744 & 13 & 240028.196 \\
1 & 349320.551 & 14 & 150543.659 \\
2 & 157067.547 & 15 & 275292.814 \\
3 & 135052.934 & 16 & 301101.219 \\
4 & 185318.110 & 17 & 100815.401 \\
5 & 118927.980 & 18 & 123254.119 \\
\textbf{6} & \textbf{921.545} & 19 & 73448.724 \\
7 & 290533.928 & 20 & 508237.438 \\
8 & 167609.763 & 21 & 199549.637 \\
9 & 357275.877 & 22 & 117261.345 \\
10 & 144992.519 & 23 & 220336.935 \\
11 & 564891.560 & 24 & 259621.119 \\
12 & 68513.462 & 25 & 176043.833 \\
\bottomrule
\end{tabular}
\end{center}

$k=6$의 값이 압도적으로 작다.
문제에서 주어진 암호문을 $k=6$ (혹은 알파벳 G)을 사용하여 복호화 시

THE SOVEREIGNTY OF THE REPUBLIC OF KOREA SHALL RESIDE IN THE 
PEOPLE AND ALL STATE AUTHORITY SHALL EMANATE FROM THE PEOPLE

로 대한민국 헌법 영문본을 구할 수 있다.

첨부된 \texttt{classifier.py --full-caesar-table}은 각 후보의
A--Z 횟수까지 출력하며, \texttt{--dump-caesar-dir DIR}은 26개 전체
복호문과 TSV를 생성한다.

\section{Vigenère key와 평문 복원}

\texttt{ciphertexts2.txt}의 4번째 줄(길이 101)에서 9글자 패턴
\texttt{OVNZPJDUO}가 인덱스 47과 92에 반복 출현한다. 간격은
$92-47=45$이고, $45$의 약수 중 조건 $r\le5$를 만족하는 값은
$\{1,3,5\}$라 이 하나만으로는 key 길이 3과 5를 구분할 수 없다.

독립된 두 번째 반복을 8번째 줄에서 찾았다: 3글자 패턴 \texttt{KWG}가
인덱스 20과 75에 반복되어 간격 $75-20=55$를 준다. 두 간격의
최대공약수
\[
 \gcd(45,55)=5
\]
가 $r\le5$ 조건 아래 유일하게 남는 값이므로 key 길이를 5로 특정한다.


이는 5-gram 통계로도 확인이 가능하다. 

원본 두 파일은 157개 줄이고 대응 줄 길이가 같다. 반복 5-gram의
인접 출현 거리를 줄 경계를 넘지 않게 수집한 결과는 다음과 같다.

\begin{center}
\begin{tabular}{rrr}
\toprule
후보 period & 나누어떨어지는 거리 & 비율 \\
\midrule
2 & 178/261 & 68.199\% \\
3 & 60/261 & 22.989\% \\
4 & 56/261 & 21.456\% \\
\textbf{5} & \textbf{260/261} & \textbf{99.617\%} \\
\bottomrule
\end{tabular}
\end{center}

우연한 반복 하나 때문에 전체 거리 GCD는 1이지만, 260/261이 5의
배수이고 문제의 $|k|\le5$ 조건이 있으므로 period는 5다. 줄마다
position modulo 5가 같은 문자를 합쳐 26개 Caesar shift를 비교했다.

\begin{center}
\begin{tabular}{rrrrr}
\toprule
위치 & 열 길이 & shift & key & 최소 $\chi^2$ \\
\midrule
0 & 4047 & 10 & K & 222.518 \\
1 & 4016 & 11 & L & 179.146 \\
2 & 3975 & 21 & V & 406.928 \\
3 & 3942 & 14 & O & 231.591 \\
4 & 3917 & 9 & J & 165.249 \\
\bottomrule
\end{tabular}
\end{center}

따라서 key는 \texttt{KLVOJ}다. 이 key를 각 줄 처음부터 다시 적용해야
한다는 사실도 두 방식의 문자별 대조로 검증했다.

\begin{center}
\begin{tabular}{lr}
\toprule
key 적용 방식 & Caesar 평문과 일치 \\
\midrule
줄마다 position 초기화 & 19,897 / 19,897 \\
줄을 이은 연속 key & 3,076 / 19,897 \\
\bottomrule
\end{tabular}
\end{center}

키값을 적용했을 때 구한 평문은 시저 암호 복호화로 구했던 평문과 같다. 

THE SOVEREIGNTY OF THE REPUBLIC OF KOREA SHALL RESIDE IN THE 
PEOPLE AND ALL STATE AUTHORITY SHALL EMANATE FROM THE PEOPLE


\section{학습 분류기}

\subsection{데이터 분할과 feature}

50자 이상인 대응 줄 쌍을 하나의 단위로 묶었다. 원래 줄 index가
5의 배수인 쌍은 holdout, 나머지는 train으로 보내 같은 평문 Caesar
버전이 train에 있고 Vigenère 버전이 holdout에 있는 직접 누출을
막았다. 최종 표본은 train 218개, holdout 58개이며 class가 균형이다.

각 줄에서 전체 IC, Shannon entropy, 최적 Caesar 정규화 $\chi^2$,
최빈 문자 비율, lag 1--10 일치율 10개, period 2--5 열의 가중 평균
IC 4개로 총 18개 feature를 만든다. 이 feature들은 전역 shift에
불변이라 특정 Caesar shift 6은 제거된다. train 평균/표준편차로
표준화한 뒤 L2 logistic regression을 외부 패키지 없이 full-batch
gradient descent로 1,500 epoch 학습했다.

\begin{center}
\begin{tabular}{lrr}
\toprule
 & predicted Caesar & predicted Vigenère \\
\midrule
actual Caesar & 29 & 0 \\
actual Vigenère & 0 & 29 \\
\bottomrule
\end{tabular}
\end{center}

정확도 $58/58$은 \emph{동일한 대한민국 헌법 corpus와 동일한 두
암호화 key 안의 paired line holdout} 결과다. 새로운 corpus, key,
Vigenère 길이에 대한 일반화 성능으로 해석할 수 없다. 전역 shift는
제거했지만 고정 period 5와 corpus에 특화될 가능성은 남아 있다.
배포물에 외부 corpus가 없으므로 별도 corpus/key 실험 결과를
만들어 적지 않았다.

\subsection{네 짧은 암호문과 최종 라벨}

\begin{center}
\small
\begin{tabular}{rrrrl}
\toprule
번호 & $P(V)$ & 모델 구조 판정 & 최적 shift & 출제 의도 추정 \\
\midrule
1 & 0.017058 & Caesar-like & 6 & Caesar \\
2 & 0.017058 & Caesar-like & 10 & Vigenère \\
3 & 0.089193 & Caesar-like & 10 & Vigenère \\
4 & 0.089193 & Caesar-like & 6 & Caesar \\
\bottomrule
\end{tabular}
\end{center}

네 암호문은 모두 하나의 고정 shift로 자연스러운 영어가 되어 관측
구조는 period-1이다.

2번을 Vigenère key \texttt{K}로 만들었다는 가설과 Caesar shift
10으로 만들었다는 가설은 암호문만으로 구별할 수 없다. 다만 shift
6이 복원된 Caesar key, shift 10이 \texttt{KLVOJ}의
첫 글자 K와 일치한다. 이 생성기 단서에 따른 출제 의도상 직접 답은
\[
 \boxed{\text{Caesar,\ Vigenère,\ Vigenère,\ Caesar}}
\]
로 제출한다. 동시에 이 목록은 식별 가능한 암호 구조가 아니라 숨은
생성 방식을 추정한 값임을 명시한다.

\section{재현}

Python 3.10 이상 표준 라이브러리만 필요하다.

\begin{lstlisting}
python3 classifier.py --problem-zip 1_암호분석.zip
python3 classifier.py --problem-zip 1_암호분석.zip --full-caesar-table
python3 classifier.py --problem-zip 1_암호분석.zip \
  --dump-caesar-dir /tmp/problem01-caesar
\end{lstlisting}

실행은 다음을 출력하고 검증한다.

\begin{itemize}
  \item IC, 모든 Caesar 후보, Kasiski와 열별 key
  \item line-reset 대조와 paired holdout confusion matrix
  \item 네 표본의 확률, shift, 평문
\end{itemize}

두 번째 명령은 26개 key 각각의 복호문과 A--Z 빈도, 카이제곱 값을
표준 출력에 모두 싣는다. 세 번째 명령은 같은 26개 복호문과 빈도 TSV를
개별 파일로 저장한다. 따라서 본문의 요약 순위표뿐 아니라 원문 3-가가
요구한 모든 key별 결과를 제출 코드에서 그대로 재생성할 수 있다.

마지막 \texttt{self-checks: PASS}가 모든 수치와 답을 검증한다.

\section{참고 자료}

\begin{itemize}
  \item William F. Friedman,
    \href{https://www.nsa.gov/Portals/75/documents/news-features/declassified-documents/friedman-documents/publications/FOLDER_233/41761039080018.pdf}
    {\emph{The Index of Coincidence and Its Applications in Cryptanalysis}}:
    단일·다중 알파벳 암호의 빈도와 IC 판정 근거.
  \item Friedrich W. Kasiski,
    \href{https://books.google.de/books?id=fB5dAAAAcAAJ}
    {\emph{Die Geheimschriften und die Dechiffrir-Kunst}}:
    반복 문자열 간격의 공약수로 key 길이를 추정하는 절차.
  \item Karl Pearson,
    \href{https://doi.org/10.1080/14786440009463897}
    {\emph{On the Criterion that a Given System of Deviations from the
    Probable \ldots}}: 관측 빈도와 영어 기대빈도의 카이제곱 비교.
  \item Claude E. Shannon,
    \href{https://doi.org/10.1002/j.1538-7305.1948.tb01338.x}
    {\emph{A Mathematical Theory of Communication}}:
    분류기의 문자분포 entropy feature.
  \item D. R. Cox,
    \href{https://doi.org/10.1111/j.2517-6161.1958.tb00292.x}
    {\emph{The Regression Analysis of Binary Sequences}}:
    binary logistic regression 모델의 배경.
\end{itemize}

\end{document}
